\section{Theoretical Background}
\label{sec:fund}

\subsection{Neural Radiance Fields and Volumetric Rendering}
The foundational architecture of NeRFs utilizes Multilayer Perceptrons (MLPs) to approximate a continuous 5D scene representation \cite{mildenhall2020nerfrepresentingscenesneural}. This function takes a spatial position $\mathbf{x} = (x, y, z)$ and a viewing direction $\mathbf{d} = (\theta, \phi)$ as inputs, outputting the volumetric density $\sigma(\mathbf{x})$ and associated color $\mathbf{c} = (r, g, b)$:

$$
\gamma(\mathbf{x}, \mathbf{d}) \rightarrow \mathbf{v}, \quad F_\theta'(\mathbf{v}) \rightarrow (\sigma, \mathbf{c}), \quad \sigma \in \mathbb{R}_{\geq 0},\ \mathbf{c} \in [0, 1]^3
$$
where $\theta$ represents the neural network weights, $F_\theta'$ is the original MLP, and $\gamma$ is a positional encoder mapping inputs into a higher-dimensional feature vector $\mathbf{v}$.

To synthesize novel views, the expected color of a pixel is computed via volumetric rendering. For a ray $\mathbf{r}(t) = \mathbf{o} + t\mathbf{d}$ with depth bounds $(t_n, t_f)$, the resulting color is:
$$\hat{C}(\mathbf{r}) = \int_{t_n}^{t_f} T(t) \sigma(\mathbf{r}(t)) \mathbf{c}(\mathbf{r}(t), \mathbf{d}) \, dt$$
where $T(t) = \exp\left(-\int_{t_n}^t \sigma(\mathbf{r}(s)) \, ds\right)$ denotes the accumulated transmittance. In practice, this integral is numerically approximated via alpha compositing over discrete samples along each ray.

\subsection{Architectural Optimizations and Generalization}
Although the original NeRF achieves high-fidelity results, its reliance on large MLPs imposes computational bottlenecks. To accelerate optimization, hybrid representations were introduced to replace Fourier positional encoding with hash encoding. Instant-NGP proposes a multiscale hash table that maps spatial coordinates to feature vectors, heavily reducing optimization time \cite{muller2022}. Nerfacto builds upon this mechanism by combining multiresolution hash encoding with proposal sampling, pose refinement, spherical harmonics, and appearance embeddings to balance visual quality and rendering speed \cite{Tancik2023}.

These per-scene methods, however, fundamentally rely on dense multi-view observations. In practical scenarios, such as mobile scanning, virtual try-on, and medical imaging (for instance, histology and holotomography), multi-view acquisition is often physically limited \cite{ijms24076747}. Under these sparse-view constraints, lacking camera pose redundancy, per-scene models suffer geometric degradation.

To synthesize novel views from sparse inputs, few-shot generalizable NeRFs incorporate structural priors learned from large datasets via dedicated feature extractors. PixelNeRF  applies a pre-trained Residual ResNet to extract spatial features from 2D projections, conditioning the radiance MLP \cite{yu2021pixelnerfneuralradiancefields}. GNT replaces the CNN with a vision transformer for global feature extraction, and substitutes the MLP with a cross-attention transformer to aggregate multi-view information \cite{t2023attentionnerfneeds}.

To maximize visual fidelity, generalizable models frequently employ fine-tuning (FT). Instead of purely zero-shot (ZS) inference, FT updates the network weights using the sparse views of the novel scene as direct photometric supervision. This rapid optimization reduces the domain gap, though it carries the risk of representational collapse if overfitted.


Given the fundamental architectural differences between per-scene optimizations and generalizable priors, selecting the optimal model for sparse-view deployments remains non-trivial. Consequently, a standardized evaluation framework is essential to systematically benchmark these trade-offs, which forms the core methodology of this work.
